% =====================================================================
%  Webseite Förderverein KF/JF Kreuzau -- Betriebshandbuch für die Betreuung
%  Vorlage: deb-support (debsupport.sty, siehe latex/README.md), Vereinskopf
%  Grundlage: deploy/setup-server-fv-web.sh, README.md (Stand siehe Titel)
%  Übersetzen mit:  pdflatex betriebshandbuch_webseite.tex  (3x)
% =====================================================================
\documentclass[a4paper,fontsize=10pt]{scrartcl}
\usepackage[bericht]{debsupport}
\tcbuselibrary{listings}
\usepackage{tikz}
\usetikzlibrary{positioning,arrows.meta,fit,backgrounds}
\setcounter{tocdepth}{2}

% Vereinskopf statt deb-support (Vorlage bleibt, nur die Absenderdaten wechseln)
\dssetup{firma=Förderverein KF/JF Kreuzau}
\makeatletter
\renewcommand*{\ds@kopf}{%
  \put(24.5,279.4){\makebox(0,0)[lb]{\dslogo[20]}}%
  \put(195.6,276.3){\makebox(0,0)[rb]{{\dsserif\fontsize{12}{14.45}\selectfont
     \begin{tabular}[b]{@{}r@{}}Kinder- und Jugendfeuerwehr\\Kreuzau\\
     foerderverein-kf-jf-kreuzau.de\end{tabular}}}}%
  \put(17.8,271.85){\rule{179.8mm}{0.4pt}}}
\renewcommand*{\dsfolgekopf}[1]{%
  \parbox[b]{\textwidth}{\dslogo[11]\hfill #1\\[-0.4ex]%
    \rule{\textwidth}{0.4pt}}}
\makeatother
\renewcommand*{\dsdokumentart}{Betriebshandbuch}
\renewcommand*{\dskurztitel}{Webseite und Mail -- Betriebshandbuch}
\renewcommand*{\dsbeleg}{}
\setlength{\emergencystretch}{3em}

% Befehle: grau hinterlegt, Umlaute erlaubt
\newtcblisting{befehl}{listing only, breakable, enhanced, colback=dshell,
  boxrule=0pt, arc=0pt, left=2mm, right=2mm, top=1mm, bottom=1mm,
  listing options={basicstyle=\small\ttfamily, breaklines=true, columns=fullflexible,
    keepspaces=true, upquote=true,
    literate={ä}{{\"a}}1 {ö}{{\"o}}1 {ü}{{\"u}}1 {Ä}{{\"A}}1 {Ö}{{\"O}}1
             {Ü}{{\"U}}1 {ß}{{\ss}}1 {–}{{--}}1 {…}{{\dots}}1 {→}{{$\rightarrow$}}1}}
\newcommand*{\cmd}[1]{\texttt{#1}}
\newcommand*{\pfad}[1]{\texttt{#1}}
% Abhaken bei der Einrichtung
\newcommand*{\erledigt}{\makebox[1.4em][l]{\setlength{\fboxsep}{0pt}\fbox{\rule{0pt}{0.7em}\hspace{0.7em}}}}

\begin{document}

\dstitel{Betriebshandbuch der Webseite und des Mailservers}
        {Für die Betreuung: Server einrichten, Mailserver, Updates, Einstellungen,
         Datensicherung, Störungen, Wartung}
        {Webseite: & \url{https://foerderverein-kf-jf-kreuzau.de}\\
         Quellcode: & \url{https://github.com/brami/FreiVerein}\\
         Server:   & vServer mit Debian 13, Admin-Benutzer \cmd{fkjs}\\
         Stand:    & \today}

\begin{dshinweis}[title={Die wichtigsten Befehle (auf dem Server, als fkjs)}]
\begin{tabularx}{\linewidth}{@{}>{\ttfamily}lX@{}}
sudo fv-web-update            & neue Version von GitHub (\cmd{main}) installieren \\
sudo fv-web-update -{}-zurueck & vorige Version wieder aktivieren \\
sudo fv-web-update -{}-status  & installierte Versionen anzeigen \\
sudo fv-web-backup            & Datensicherung sofort ausführen \\
sudo fv-web-restore -t        & Wiederherstellungstest (ändert nichts) \\
sudo fv-manage <befehl>       & Befehle der Anwendung (siehe Abschnitt~\ref{sec:verwaltung}) \\
systemctl status fv-web       & läuft die Webseite? \\
sudo fv-web-status            & Gesamtzustand (wie \cmd{/health}) \\
journalctl -u fv-web -e       & Protokoll der Webseite \\
sudo fv-mail status           & Mailserver: Dienste, Warteschlange, Zertifikat \\
sudo fv-mail dns              & Mailserver: DNS-Einträge anzeigen und prüfen \\
\end{tabularx}
\end{dshinweis}

\tableofcontents

% ==================================================================
\section{Überblick}

\subsection{Aufbau}

Die Webseite ist eine eigene Anwendung auf Basis von Django. Sie läuft allein auf
einem vServer, ausschließlich mit Paketen aus Debian (kein Docker). Den Quellcode
holt der Server selbst von GitHub. Auf demselben Server läuft der Mailserver des
Vereins; Adressen, Postfächer und Verteiler pflegt der Vorstand in der Webseite.

\begin{figure}[h]
\centering
\begin{tikzpicture}[font=\small,
    box/.style={draw, rounded corners=3pt, align=center, minimum height=1cm,
                text width=3.1cm, fill=white},
    daten/.style={draw, align=center, text width=3.1cm, minimum height=0.9cm, fill=dshell},
    pf/.style={-{Stealth[length=2.3mm]}, thick}, lbl/.style={font=\footnotesize}]
  \node[box] (browser) {\textbf{Besucher}\\Browser};
  \node[box, right=1.1cm of browser] (caddy) {\textbf{Caddy}\\HTTPS, Zertifikat,\\statische Dateien};
  \node[box, right=1.1cm of caddy, draw=dsrot, very thick] (app) {\textbf{fv-web}\\gunicorn + Django};
  \node[daten, below=0.9cm of app] (db) {\textbf{PostgreSQL 17}\\Datenbank fvweb};
  \node[daten, below=0.9cm of caddy] (files) {\textbf{Anhänge}\\/var/lib/fv-web/private};
  \node[box, right=1.1cm of db] (mail) {\textbf{Mailserver}\\Exim, Dovecot, rspamd};
  \node[box, above=0.9cm of mail] (prog) {\textbf{Mailprogramme}\\IMAP, SMTP};
  \node[box, below=0.9cm of files] (backup) {\textbf{fv-web-backup}\\täglich 03:30 (restic)};
  \node[daten, left=1.1cm of backup] (ziel) {\textbf{Backup-Ziel}\\lokal oder extern};
  \node[daten, below=0.9cm of mail] (vmail) {\textbf{Postfächer}\\/var/vmail};
  \draw[pf, <->] (browser) -- (caddy);
  \draw[pf, <->] (caddy) -- node[lbl, above]{Socket} (app);
  \draw[pf, <->] (app) -- (db);
  \draw[pf, <->] (app) -- (files);
  \draw[pf] (mail) -- node[lbl, above]{liest} (db);
  \draw[pf] (app) -- node[lbl, above, sloped]{SMTP} (mail);
  \draw[pf, <->] (prog) -- (mail);
  \draw[pf, <->] (mail) -- (vmail);
  \draw[pf] (db) |- (backup);
  \draw[pf] (files) -- (backup);
  \draw[pf] (vmail) |- (backup);
  \draw[pf] (backup) -- (ziel);
\end{tikzpicture}
\caption{Bestandteile auf dem Server}
\end{figure}

\subsection{Dienste}

\begin{center}
\small
\begin{tabularx}{\textwidth}{@{}>{\ttfamily}lX@{}}
\toprule
\textbf{\sffamily Dienst} & \textbf{Aufgabe} \\
\midrule
fv-web.service            & die Webseite (gunicorn, 3 Prozesse, Socket \pfad{/run/fv-web/gunicorn.sock}) \\
fv-web-newsletter.service & verschickt Newsletter; startet, stoppt und startet neu zusammen mit fv-web \\
caddy.service             & HTTPS, Zertifikat von Let's Encrypt (automatisch), statische Dateien \\
postgresql.service        & Datenbank \\
fv-web-backup.timer       & tägliche Sicherung gegen 03:30 \\
fv-web-status.timer       & alle 5 Minuten: Zustand für \cmd{/health} prüfen \\
unattended-upgrades       & automatische Sicherheitsupdates von Debian \\
fail2ban, ufw             & Schutz gegen falsche Passwörter (SSH, Mail), Firewall \\
\midrule
exim4.service             & Mailserver: empfängt (Port 25) und verschickt Mails, nimmt Mails der
                            Mailprogramme an (465, 587) \\
dovecot.service           & Postfächer: IMAP (993), eigene Filterregeln (4190), Speichergrenze \\
rspamd.service, redis     & Spam, Greylisting, unzulässige Anhänge \\
unbound.service           & eigener DNS-Resolver mit DNSSEC-Prüfung \\
fv-mail-sync.timer        & alle 5 Minuten: Abwesenheitsnotizen, Archiv-Rechte, Belegung,
                            Plus-Adressen mit der Webseite abgleichen \\
fv-mail-taeglich.timer    & täglich 04:15: Erinnerungen und Löschungen gesperrter Postfächer \\
certbot.timer             & verlängert das Zertifikat für \cmd{mail.<domain>} \\
\bottomrule
\end{tabularx}
\end{center}

\subsection{Verzeichnisse und Dateien}

\begin{center}
\small
\begin{tabularx}{\textwidth}{@{}>{\ttfamily}lX@{}}
\toprule
\textbf{\sffamily Pfad} & \textbf{Inhalt} \\
\midrule
/srv/fv-web/releases/      & installierte Versionen (die letzten drei) \\
/srv/fv-web/current        & Verweis auf die aktive Version \\
/etc/fv-web/env            & Einstellungen der Anwendung (Abschnitt~\ref{sec:einstellungen}) \\
/etc/fv-web/deploy.env     & Repository, Standard-Branch, Domain (nur root) \\
/etc/fv-web/backup.env     & Datensicherung inkl. \textbf{Passwort} (nur root) \\
/var/lib/fv-web/private    & Anhänge und Bilder (nur für die Anwendung lesbar) \\
/var/lib/fv-web/static     & CSS, JavaScript, Editor (liefert Caddy aus) \\
/var/lib/fv-web/mails      & E-Mails als Dateien, solange kein Mailserver eingetragen ist \\
/var/lib/fv-web/wartung    & Wartungsseite (schreibt \cmd{fv-web-update}) \\
/var/backups/fv-web        & Datenbank-Dumps der letzten 7 Tage \\
/var/log/fv-web-update.log & Protokoll aller Updates \\
/var/log/caddy/fv-web.log  & Zugriffe auf die Webseite \\
/etc/caddy/Caddyfile       & wird von \cmd{fv-web-update} geschrieben -- nicht von Hand ändern \\
\midrule
/var/vmail/konten/<nr>     & Postfach des Kontos Nummer \cmd{<nr>} (Maildir) \\
/var/vmail/archiv          & Archive der Verteiler \\
/etc/fv-mail/mail.env      & Mailserver: Rechnername, DKIM-Selektor, Geheimnisse (nur root) \\
/etc/fv-mail/dkim, /tls    & DKIM-Schlüssel, Zertifikat für \cmd{mail.<domain>} \\
/etc/fv-mail/aliases       & Ziel für postmaster@, abuse@, root@ (die Betreuung) \\
/etc/exim4/exim4.conf, /etc/dovecot/dovecot.conf, /etc/rspamd/local.d/
                           & schreibt das Einrichtungsskript -- nicht von Hand ändern \\
/var/log/exim4/mainlog     & alle Mails (ohne Betreff), 14 Tage \\
\bottomrule
\end{tabularx}
\end{center}

\subsection{Konten}

\begin{description}[style=nextline,leftmargin=0pt,itemsep=0.3\baselineskip]
  \item[\cmd{fkjs} (Linux)] Admin-Benutzer des Servers; Anmeldung nur per SSH-Schlüssel,
    \cmd{sudo} ohne Passwort. Anmeldung als root und mit Passwort ist abgeschaltet.
  \item[\cmd{fvweb} (Linux und PostgreSQL)] Dienstbenutzer der Anwendung; Datenbankzugriff
    nur lokal über den Socket, ohne Passwort.
  \item[\cmd{dovecot}, \cmd{Debian-exim} (PostgreSQL)] Rollen des Mailservers; dürfen nur
    die für sie bestimmten Sichten lesen (Adressen, Verteiler, Passwort-Hashes), nichts
    schreiben. Die Leserechte vergibt \cmd{fv-web-update} bei jeder Installation neu.
  \item[\cmd{vmail} (Linux)] Besitzer aller Postfächer.
  \item[\cmd{fv-archiv@<domain>}] internes Postfach für die Archive der Verteiler; keine
    Anmeldung möglich.
  \item[Administrator der Webseite] Konto der Webseite mit vollen Rechten
    (\emph{Superuser}); sieht zusätzlich \emph{Stufen} und \emph{Einstellungen der
    Webseite}. Anlegen siehe Abschnitt~\ref{sec:verwaltung}.
\end{description}

% ==================================================================
\section{Server einrichten}\label{sec:einrichten}

Das Skript \pfad{deploy/setup-server-fv-web.sh} läuft auf dem eigenen Rechner und
richtet den Server per SSH ein. Es kann beliebig oft erneut laufen; jeder Schritt
prüft, ob er schon erledigt ist. \cmd{-h} zeigt alle Optionen.

\subsection{Einstellungsdatei}\label{sec:einstellungsdatei}

Alle dauerhaften Einstellungen eines Servers -- Server, Domain, www, Mailserver,
Betreuungsadresse, Backup-Ziel, Admin-Benutzer, SSH-Port -- stehen in einer
\textbf{Einstellungsdatei} auf dem eigenen Rechner, z.\,B.\ \pfad{kreuzau.env}. Das
Skript liest sie mit \cmd{-f}, übernimmt Änderungen aus den übrigen Optionen und
schreibt sie zurück; geänderte Werte zeigt es an. Spätere Aufrufe brauchen dann nur
noch \cmd{-f} und das, was diesmal anders sein soll. So geht keine Einstellung
verloren, und jeder Aufruf richtet den Server gleich ein.

\begin{befehl}
SERVER="203.0.113.10"
DOMAIN="foerderverein-kf-jf-kreuzau.de"
WITH_WWW="1"
WITH_MAIL="1"
MAIL_HOST=""                       (leer: mail.<domain>)
POSTMASTER="betreuung@example.org"
BACKUP_REPO="sftp://u12345@u12345.your-storagebox.de:23/fv-web"
…
\end{befehl}

\begin{itemize}
  \item Gespeichert werden \cmd{-H -d -A -W -M -X -P -e -B -R -u -n -t -p -S}. Einmalige
        Optionen (\cmd{-c -U -N -L -r -a -m -y}) nicht.
  \item Etwas abschalten (z.\,B.\ www): in der Datei auf \cmd{0} setzen.
  \item Die Datei enthält keine Passwörter. Sie gehört zu den Unterlagen der Betreuung
        (Passwortsafe oder Projektablage), nicht ins öffentliche Repository.
  \item Auch ohne Datei entfernt das Skript ein eingerichtetes www nicht, nur weil
        \cmd{-W} fehlt.
\end{itemize}

\subsection{Voraussetzungen}

\begin{itemize}
  \item vServer mit frischem Debian 13, Anmeldung als root mit SSH-Schlüssel.
  \item DNS: A-Eintrag (und ggf.\ AAAA) der Domain zeigt auf den Server; für
        \cmd{-W} auch \cmd{www}.
  \item Auf dem eigenen Rechner: \cmd{bash}, \cmd{ssh}; das Skript allein genügt
        (das Repository ist nur für \cmd{-L} nötig).
\end{itemize}

\subsection{Ablauf}

\begin{enumerate}
  \item \textbf{Nur prüfen} (ändert nichts): SSH, belegte Ports 80/443, DNS.
\begin{befehl}
./setup-server-fv-web.sh -f kreuzau.env -H 203.0.113.10 -d foerderverein-kf-jf-kreuzau.de -W -c
\end{befehl}
  \item \textbf{Erstinstallation} mit Administrator-Konto der Webseite und externem
        Backup-Ziel; legt \pfad{kreuzau.env} an:
\begin{befehl}
./setup-server-fv-web.sh -f kreuzau.env -H 203.0.113.10 -d foerderverein-kf-jf-kreuzau.de -W \
    -n fv-web -a michael.bramer -m adresse@example.org \
    -B sftp://u12345@u12345.your-storagebox.de:23/fv-web
\end{befehl}
        Das Skript aktualisiert das System, legt \cmd{fkjs} an (übernimmt die
        SSH-Schlüssel von root), richtet Firewall, fail2ban, PostgreSQL, Caddy,
        Dienste und Sicherung ein, testet die Anmeldung als \cmd{fkjs} und härtet erst
        danach SSH. Dann holt der Server die Anwendung von GitHub und installiert sie.
  \item \textbf{Ausgaben beachten:}
    \begin{itemize}
      \item den \textbf{Link zum Setzen des Passworts} für das Administrator-Konto
            (7 Tage gültig),
      \item den \textbf{öffentlichen Backup-Schlüssel}: bei der Storage-Box als
            erlaubten Schlüssel eintragen, sonst schlägt die Sicherung fehl,
      \item bei privatem Repository den \textbf{Deploy-Schlüssel}: bei GitHub unter
            \emph{Settings $\rightarrow$ Deploy keys} eintragen (ohne Schreibrecht),
            dann erneut mit \cmd{-U} aufrufen.
    \end{itemize}
  \item \textbf{Passwort der Datensicherung sichern} -- siehe Hinweis in
        Abschnitt~\ref{sec:sicherung}.
  \item Mit dem Link Passwort setzen, anmelden, Zwei-Faktor-Anmeldung einrichten.
  \item In der Verwaltung unter \emph{Einstellungen der Webseite} die \emph{Weitere
        Antwortadresse} (Vorstand) eintragen; \emph{Stufen} prüfen.
  \item \textbf{Mailserver} einrichten: Abschnitt~\ref{sec:mail}.
  \item Vorstandsmitglieder in der Verwaltung anlegen und einladen (siehe
        Vorstandsanleitung). Solange kein Mailserver eingetragen ist, zeigt die
        Verwaltung den Einladungslink an; er muss dann selbst weitergegeben werden.
\end{enumerate}

Protokolle der Einrichtung: lokal \pfad{logs/setup-<server>-<zeit>.log} neben dem
Skript, auf dem Server \pfad{/var/log/fv-web-setup/}.

% ==================================================================
\section{Mailserver einrichten -- Schritt für Schritt}\label{sec:mail}

Den Mailserver richtet das Einrichtungsskript mit \cmd{-M} ein. Einiges lässt sich
aber nur \textbf{von Hand} erledigen: Einträge im DNS (bei SchlundTech), der
Reverse-DNS-Eintrag beim Server-Anbieter und die Tests. Diese Schritte stehen hier
in der richtigen Reihenfolge, zum Abhaken. In den Beispielen steht
\cmd{foerderverein-kf-jf-kreuzau.de} für die Domain, \cmd{203.0.113.10} und
\cmd{2001:db8::10} für die Adressen des Servers.

\begin{dshinweis}[title={Reihenfolge ist wichtig}]
Der \textbf{MX-Eintrag} (Schritt~\ref{mx}) wird als \emph{letzter} umgestellt. Bis dahin
kommen Mails von außen weiter beim bisherigen Anbieter an, und alles andere lässt sich
in Ruhe einrichten und testen. Gibt es noch keinen bisherigen Anbieter, kann der MX
gleich mit den übrigen Einträgen gesetzt werden.
\end{dshinweis}

\subsection{Vorbereitung (einige Tage vorher)}

\begin{enumerate}
  \item\label{ttl} \erledigt\textbf{TTL senken.} Bei SchlundTech in der Zone der Domain die
        TTL des bestehenden MX-Eintrags auf \cmd{300} (5~Minuten) setzen. Danach
        mindestens so lange warten wie die \emph{alte} TTL (meist einen Tag). Dann
        wirkt die spätere Umstellung schnell, und ein Zurück ist schnell möglich.
  \item \erledigt\textbf{Port 25 beim Server-Anbieter.} Manche Anbieter sperren ausgehende
        Verbindungen auf Port~25 bei neuen Servern. Das Skript prüft das und warnt;
        im Zweifel vorher im Kundenmenü nachsehen oder beim Support freischalten lassen.
  \item \erledigt\textbf{Reverse-DNS (PTR) beim Server-Anbieter} setzen, im Kundenmenü
        des Servers (nicht bei SchlundTech):
\begin{befehl}
203.0.113.10   →  mail.foerderverein-kf-jf-kreuzau.de
2001:db8::10   →  mail.foerderverein-kf-jf-kreuzau.de   (falls der Server IPv6 hat)
\end{befehl}
        Ohne passenden PTR landen Mails bei großen Anbietern im Spam oder werden
        abgewiesen. Lässt sich für IPv6 kein PTR setzen, verschickt der Mailserver
        automatisch nur über IPv4.
  \item \erledigt\textbf{Adresse der Betreuung festlegen}, an die postmaster@, abuse@ und
        root@ gehen (eine Adresse \emph{außerhalb} der eigenen Domain, damit Meldungen
        auch bei einer Störung des Mailservers ankommen).
\end{enumerate}

\subsection{DNS für den Mailserver-Namen}

\begin{enumerate}[resume]
  \item \erledigt\textbf{Neue Einträge bei SchlundTech} (\emph{Domains $\rightarrow$
        DNS-Einstellungen bzw. Zone bearbeiten}):
\begin{befehl}
mail.foerderverein-kf-jf-kreuzau.de.     A      203.0.113.10
mail.foerderverein-kf-jf-kreuzau.de.     AAAA   2001:db8::10     (nur mit IPv6 und PTR)
mta-sts.foerderverein-kf-jf-kreuzau.de.  A      203.0.113.10
\end{befehl}
        Die bestehenden Einträge der Webseite und den MX \textbf{noch nicht} ändern.
  \item \erledigt\textbf{Prüfen}, ob die Einträge sichtbar sind (auf dem eigenen Rechner):
\begin{befehl}
dig +short A mail.foerderverein-kf-jf-kreuzau.de @1.1.1.1
dig +short -x 203.0.113.10 @1.1.1.1
\end{befehl}
        Beide müssen das Erwartete zeigen, bevor es weitergeht.
\end{enumerate}

\subsection{Einrichtungsskript}

\begin{enumerate}[resume]
  \item \erledigt\textbf{Skript mit \cmd{-M} aufrufen.} Auf einem schon eingerichteten Server
        einmal \emph{ohne} \cmd{-U}, damit auch die Datensicherung die Postfächer
        aufnimmt. \cmd{-M} und \cmd{-P} landen in der Einstellungsdatei; danach genügt
        \cmd{-f kreuzau.env -U}.
\begin{befehl}
./setup-server-fv-web.sh -f kreuzau.env -M -P betreuung@example.org
\end{befehl}
        Das Skript installiert Exim, Dovecot, rspamd, Redis, unbound und certbot,
        holt das Zertifikat für \cmd{mail.<domain>}, erzeugt den DKIM-Schlüssel, öffnet
        die Ports 25, 465, 587, 993 und 4190, stellt die Webseite auf den eigenen
        Mailserver um und gibt am Ende die nötigen DNS-Einträge aus.
  \item \erledigt\textbf{Ausgaben beachten:}
    \begin{itemize}
      \item \emph{WARNUNG: PTR \dots} -- Reverse-DNS stimmt noch nicht (Schritt~3).
      \item \emph{WARNUNG: Ausgehender Port 25 scheint gesperrt} -- Anbieter fragen.
      \item \emph{HINWEIS: \dots\ zeigt noch nicht auf \dots\ – kein Zertifikat} -- dann
            gilt vorläufig ein selbst signiertes Zertifikat; nach dem DNS-Eintrag das
            Skript noch einmal mit \cmd{-f kreuzau.env -U} aufrufen.
      \item \emph{FEHLER: Dovecot-Konfiguration ungültig} -- die bisherige bleibt aktiv;
            Meldung notieren und an die Entwicklung geben (\cmd{sudo doveconf -n -c
            /etc/dovecot/dovecot.conf.neu}).
    \end{itemize}
\end{enumerate}

\subsection{DNS-Einträge für Mail}

\begin{enumerate}[resume]
  \item \erledigt\textbf{Einträge anzeigen} (auf dem Server):
\begin{befehl}
sudo fv-mail dns
\end{befehl}
        Die Ausgabe enthält alle Einträge mit den richtigen Werten, insbesondere den
        DKIM-Schlüssel. Sie sieht etwa so aus:
\begin{befehl}
foerderverein-kf-jf-kreuzau.de.             TXT  "v=spf1 mx -all"
mail.foerderverein-kf-jf-kreuzau.de.        TXT  "v=spf1 a -all"
fv2026._domainkey.foerderverein-kf-jf-kreuzau.de.  TXT  "v=DKIM1; k=rsa; p=MIIBIjAN…"
_dmarc.foerderverein-kf-jf-kreuzau.de.      TXT  "v=DMARC1; p=none; rua=mailto:postmaster@…"
_mta-sts.foerderverein-kf-jf-kreuzau.de.    TXT  "v=STSv1; id=20261113120000"
_smtp._tls.foerderverein-kf-jf-kreuzau.de.  TXT  "v=TLSRPTv1; rua=mailto:postmaster@…"
_submissions._tcp.foerderverein-kf-jf-kreuzau.de.  SRV  0 1 465 mail.foerderverein-kf-jf-kreuzau.de.
_imaps._tcp.foerderverein-kf-jf-kreuzau.de.        SRV  0 1 993 mail.foerderverein-kf-jf-kreuzau.de.
\end{befehl}
  \item \erledigt\textbf{Bei SchlundTech eintragen} -- alle TXT- und SRV-Einträge aus der
        Ausgabe, \textbf{außer dem MX}. Hinweise:
    \begin{itemize}
      \item \textbf{SPF:} Gibt es schon einen TXT-Eintrag \cmd{v=spf1 \dots} für die
            Domain (vom bisherigen Anbieter), diesen \emph{ersetzen}, nicht einen
            zweiten anlegen. Bis zur Umstellung des MX darf der alte Anbieter noch
            darin stehen, z.\,B.\ \cmd{v=spf1 mx a:mail.foerderverein-kf-jf-kreuzau.de
            include:<alter-anbieter> -all}; nach der Umstellung auf \cmd{v=spf1 mx -all}
            kürzen.
      \item \textbf{DKIM:} Der Wert ist länger als 255 Zeichen. Die Ausgabe teilt ihn
            in mehrere Teile in Anführungszeichen; SchlundTech nimmt den ganzen Wert
            meist auch am Stück. Name: \cmd{fv2026.\_domainkey} (Selektor laut Ausgabe).
      \item \textbf{DMARC} zunächst mit \cmd{p=none}: Es wird nur berichtet, nichts
            abgewiesen. Verschärfen nach einigen Wochen (Schritt~\ref{dmarc}).
      \item \textbf{SRV-Einträge} helfen Mailprogrammen bei der Einrichtung. Die
            Einträge mit Ziel \cmd{.} (Punkt) sagen „gibt es nicht“ (kein
            unverschlüsseltes IMAP, kein POP3) und sind optional.
      \item \textbf{TLSA} nur eintragen, wenn DNSSEC eingeschaltet ist (nächster Schritt).
    \end{itemize}
  \item \erledigt\textbf{DNSSEC (empfohlen, aber freiwillig).} Bei SchlundTech für die
        Domain DNSSEC einschalten. Erst wenn \cmd{dig +dnssec
        foerderverein-kf-jf-kreuzau.de SOA} im Kopf \cmd{ad} zeigt, den
        TLSA-Eintrag aus \cmd{fv-mail dns} setzen:
\begin{befehl}
_25._tcp.mail.foerderverein-kf-jf-kreuzau.de.  TLSA  3 1 1 <wert aus fv-mail dns>
\end{befehl}
        Der Wert bleibt bei Zertifikatsverlängerungen gleich (der Schlüssel wird
        weiterverwendet). Wird der Schlüssel doch einmal ausgetauscht, \emph{zuerst} den
        neuen TLSA-Eintrag \emph{zusätzlich} setzen und erst nach einem Tag den alten
        entfernen.
  \item \erledigt\textbf{Skript noch einmal} mit \cmd{-f kreuzau.env -U} aufrufen. Jetzt richtet es
        MTA-STS ein (\cmd{https://mta-sts.<domain>}) und holt ggf.\ das Zertifikat.
  \item \erledigt\textbf{Prüfen:}
\begin{befehl}
sudo fv-mail dns       # unten: alle Zeilen "OK" außer MX (kommt später)
sudo fv-mail status    # alle Dienste "active", Zertifikat von Let's Encrypt
\end{befehl}
\end{enumerate}

\subsection{Testen vor der Umstellung}

\begin{enumerate}[resume]
  \item \erledigt\textbf{Eigenes Postfach anlegen:} in der Verwaltung beim eigenen Konto
        eine Vereinsadresse eintragen und \emph{Postfach} anhaken. Das Passwort ist das
        der Webseite (einmal neu anmelden, falls es noch im alten Format gespeichert
        ist).
  \item \erledigt\textbf{Mailprogramm einrichten:} Thunderbird findet die Einstellungen
        selbst (nur Adresse und Passwort eingeben); iPhone/Mac über das Profil unter
        \emph{Konto} auf der Webseite. Sonst von Hand: IMAP
        \cmd{mail.<domain>}, Port 993, SSL/TLS; SMTP \cmd{mail.<domain>}, Port 465,
        SSL/TLS; Benutzername = vollständige Vereinsadresse.
  \item \erledigt\textbf{Senden testen:} vom neuen Postfach an eine Adresse bei Gmail oder
        GMX schreiben. Dort \emph{Original anzeigen}: \cmd{SPF: PASS}, \cmd{DKIM: PASS},
        \cmd{DMARC: PASS}. Zusätzlich eine Mail an die Adresse von
        \url{https://www.mail-tester.com} schicken; Ziel: 9/10 oder besser.
  \item \erledigt\textbf{Webseite testen:} in der Verwaltung eine Einladung an sich selbst
        schicken. Sie muss ankommen (Absender \cmd{webmaster@<domain>}).
  \item \erledigt\textbf{Empfang intern testen:} vom eigenen Postfach an einen Verteiler
        schreiben, in dem man selbst steht; Abwesenheitsnotiz auf der Webseite
        einschalten und von einem anderen Postfach anschreiben.
\end{enumerate}

\subsection{Umstellung}

\begin{enumerate}[resume]
  \item\label{mx} \erledigt\textbf{MX umstellen} bei SchlundTech: den bisherigen MX-Eintrag
        ersetzen durch
\begin{befehl}
foerderverein-kf-jf-kreuzau.de.   MX   10 mail.foerderverein-kf-jf-kreuzau.de.
\end{befehl}
        Nach etwa der TTL (5~Minuten, siehe Schritt~\ref{ttl}) kommen neue Mails auf
        dem eigenen Server an. Den alten Anbieter noch einige Tage weiter abrufen; dort
        können vereinzelt Mails eintreffen, die vor der Umstellung unterwegs waren.
  \item \erledigt\textbf{Empfang von außen testen:} von Gmail/GMX an die eigene
        Vereinsadresse schreiben. Wird die erste Mail eines neuen Absenders um einige
        Minuten verzögert, ist das das Greylisting -- gewollt.
  \item \erledigt\textbf{Vorhandene Mails übernehmen} (falls es alte Postfächer im
        Maildir-Format gibt): Verzeichnis auf den Server kopieren und je Postfach
        einspielen. Vorhandene Mails im neuen Postfach bleiben erhalten.
\begin{befehl}
scp -r alt-maildir/ fkjs@foerderverein-kf-jf-kreuzau.de:/tmp/vera
sudo fv-mail import /tmp/vera vera.vorstand@foerderverein-kf-jf-kreuzau.de
sudo rm -r /tmp/vera
\end{befehl}
  \item \erledigt\textbf{SPF kürzen} auf \cmd{v=spf1 mx -all}, sobald der alte Anbieter
        nichts mehr für die Domain verschickt.
  \item \erledigt\textbf{Mitglieder informieren} (Anleitung für Mitglieder): neue
        Einstellungen der Mailprogramme, Passwort = Passwort der Webseite.
\end{enumerate}

\subsection{Nach zwei bis vier Wochen}

\begin{enumerate}[resume]
  \item\label{dmarc} \erledigt\textbf{DMARC-Berichte ansehen}, die an postmaster@ gehen (XML in
        ZIP-Anhängen; z.\,B.\ mit einem kostenlosen Online-Auswerter lesbar). Stammen
        alle Mails der Domain vom eigenen Server und bestehen sie SPF und DKIM, DMARC
        verschärfen:
\begin{befehl}
_dmarc.foerderverein-kf-jf-kreuzau.de.  TXT  "v=DMARC1; p=quarantine; rua=mailto:postmaster@foerderverein-kf-jf-kreuzau.de"
\end{befehl}
        Nach weiteren Wochen ohne Auffälligkeiten ggf.\ \cmd{p=reject}.
  \item \erledigt\textbf{MTA-STS scharf schalten:} auf dem Server in
        \pfad{/etc/fv-mail/mail.env} \cmd{MTA\_STS\_MODE="enforce"} setzen, dann das
        Skript mit \cmd{-f kreuzau.env -U}. Es erneuert Richtlinie und Kennung; den neuen
        \cmd{\_mta-sts}-Eintrag aus \cmd{sudo fv-mail dns} bei SchlundTech übernehmen.
  \item \erledigt\textbf{TTL zurück} auf \cmd{3600} für den MX-Eintrag.
  \item \erledigt\textbf{Reputation prüfen:} \url{https://www.spamhaus.org/lookup/} und
        \url{https://mxtoolbox.com/blacklists.aspx} für die IP-Adresse des Servers.
\end{enumerate}

\subsection{Weitere Domains (z.\,B.\ mit Umlaut oder Kurzform)}\label{sec:zusatzdomains}

Neben der Hauptdomain kann der Verein weitere Domains besitzen, z.\,B.\
\cmd{förderverein-kf-jf-kreuzau.de} oder eine Kurzform wie \cmd{fv-kf-jf-kreuzau.de}.
Jede weitere Domain
\begin{itemize}
  \item leitet die Webseite auf die Hauptdomain um (mit \cmd{-W} auch \cmd{www.}),
  \item nimmt Mails an \emph{alle} Adressen an, die es unter der Hauptdomain gibt:
        \cmd{vera@fv-kf-jf-kreuzau.de} landet im selben Postfach wie
        \cmd{vera@foerderverein-kf-jf-kreuzau.de}; ebenso Aliase, Verteiler, Sperren und
        Plus-Adressen,
  \item kann auch zum \emph{Senden} verwendet werden: Wer \cmd{vera@foerderverein-…}
        verwenden darf, darf auch \cmd{vera@fv-kf-jf-kreuzau.de} als Absender eintragen
        (im Mailprogramm als weitere Identität). Die Mails werden ebenso mit DKIM
        signiert. Üblich bleibt die Hauptdomain; bei Domains mit Umlaut wird vom Senden
        abgeraten, weil manche Empfänger solche Adressen nicht verarbeiten.
\end{itemize}

Domains mit Umlaut schreibt das Skript selbst in die technische Form (Punycode) um,
z.\,B.\ \cmd{xn--frderverein-kf-jf-kreuzau-…de}; so steht sie auch in der
Einstellungsdatei und wird so bei SchlundTech angezeigt.

\begin{enumerate}
  \item \erledigt Domain bei SchlundTech registrieren bzw.\ auf den Verein übertragen.
  \item \erledigt In ihrer Zone den A-Eintrag (und mit \cmd{-W} \cmd{www}) auf den Server
        setzen.
  \item \erledigt Skript mit der Domain aufrufen (mehrere \cmd{-A} möglich, werden in der
        Einstellungsdatei ergänzt):
\begin{befehl}
./setup-server-fv-web.sh -f kreuzau.env -U -A förderverein-kf-jf-kreuzau.de -A fv-kf-jf-kreuzau.de
\end{befehl}
  \item \erledigt \cmd{sudo fv-mail dns} zeigt je weiterer Domain die Mail-Einträge:
        \cmd{MX} auf \cmd{mail.<hauptdomain>}, \cmd{SPF "v=spf1 mx -all"}, den
        DKIM-Eintrag (derselbe Wert wie bei der Hauptdomain, gleicher Selektor) und
        \cmd{DMARC} zunächst mit \cmd{p=none}. Eintragen und mit \cmd{sudo fv-mail dns}
        prüfen. DMARC später wie bei der Hauptdomain verschärfen
        (Schritt~\ref{dmarc}).
\end{enumerate}

Entfernen: Domain in der Einstellungsdatei aus \cmd{ZUSATZ\_DOMAINS} löschen und das
Skript mit \cmd{-f kreuzau.env -U} aufrufen; danach die DNS-Einträge löschen.

\subsection{Was der Mailserver selbst erledigt}

\begin{itemize}
  \item Adressen, Aliase, Verteiler, gesperrte Adressen und Plus-Adressen liest er
        direkt aus der Datenbank -- Änderungen in der Webseite wirken sofort.
  \item Abwesenheitsnotizen, Rechte auf die Archive, die Belegung der Postfächer und
        neu benutzte Plus-Adressen gleicht \cmd{fv-mail-sync} alle 5 Minuten ab.
  \item Postfächer ausgeschiedener Personen: Erinnerung 14 Tage vor der Löschung,
        Löschung 90 Tage nach der Sperre (\cmd{fv-mail-taeglich}). Die Adresse bleibt
        danach 12 Monate gesperrt.
  \item Papierkorb und Spam-Ordner werden nach 90 Tagen geleert. Ab 90\,\% und 95\,\%
        Belegung bekommt die Person eine Hinweis-Mail; volle Postfächer nehmen nichts
        mehr an.
  \item Verschiebt jemand eine Mail in den Spam-Ordner (oder heraus), lernt rspamd
        daraus.
\end{itemize}

% ==================================================================
\section{Was automatisch läuft}

\begin{description}[style=nextline,leftmargin=0pt,itemsep=0.3\baselineskip]
  \item[Sicherheitsupdates von Debian] täglich durch \cmd{unattended-upgrades}, auch für
    PostgreSQL, Caddy und den Mailserver. Der Server startet dabei \textbf{nicht} von
    selbst neu (siehe Wartungsplan, Abschnitt~\ref{sec:wartung}).
  \item[Datensicherung] täglich gegen 03:30, sonntags mit Prüfung eines Teils der
    gesicherten Daten.
  \item[Zertifikate] Caddy holt und verlängert das Zertifikat der Webseite selbst,
    \cmd{certbot.timer} das des Mailservers.
  \item[Mailserver] Abgleich mit der Webseite alle 5 Minuten, Löschfristen täglich
    (Abschnitt~\ref{sec:mail}).
  \item[Neustart bei Absturz] Die Dienste starten nach einem Fehler automatisch neu.
\end{description}

Nicht automatisch: neue Versionen der Anwendung, Updates der Python-Pakete
(Django usw.) und Neustarts nach Kernel-Updates.

% ==================================================================
\section{Neue Version einspielen}\label{sec:update}

\subsection{Normaler Weg}

\begin{enumerate}
  \item Änderung als Pull Request auf GitHub; die automatischen Tests (SQLite,
        PostgreSQL, Installationsskript) müssen grün sein.
  \item Pull Request in \cmd{main} übernehmen.
  \item Auf dem Server einspielen:
\begin{befehl}
ssh fkjs@foerderverein-kf-jf-kreuzau.de sudo fv-web-update
\end{befehl}
\end{enumerate}

\cmd{fv-web-update} holt den Stand, baut eine eigene Python-Umgebung, prüft die
Konfiguration, \textbf{sichert die Datenbank}, migriert sie, vergibt die Leserechte
für den Mailserver, sammelt die statischen Dateien, schaltet um und prüft, ob die
Webseite antwortet (\pfad{/gesund/}). Startet die neue Version nicht, wird
automatisch die vorige wieder aktiviert. Die Webseite ist dabei nur wenige Sekunden
nicht erreichbar; in dieser Zeit (und bei jedem Ausfall der Anwendung) zeigt Caddy
eine Wartungsseite, die sich alle 30~Sekunden selbst neu lädt. Der Mailserver läuft
während des Updates weiter.

Änderungen an der Konfiguration des Mailservers (neue Fassung des
Einrichtungsskripts) werden mit \cmd{./setup-server-fv-web.sh -f kreuzau.env -U}
eingespielt.

\subsection{Weitere Möglichkeiten}

\begin{center}
\small
\begin{tabularx}{\textwidth}{@{}>{\ttfamily}lX@{}}
\toprule
sudo fv-web-update <branch|tag|commit> & bestimmten Stand installieren, z.\,B.\ zum Testen eines Branches \\
sudo fv-web-update -f                  & neu installieren, auch wenn der Stand schon aktiv ist \\
sudo fv-web-update -{}-zurueck          & vorige Version aktivieren \\
sudo fv-web-update -{}-status           & aktive und installierte Versionen \\
sudo fv-web-update -{}-pruefen          & Zugriff auf das Repository testen \\
\bottomrule
\end{tabularx}
\end{center}

Vom eigenen Rechner aus geht es auch mit dem Einrichtungsskript: \cmd{-U} spielt nur
die Anwendung ein; \cmd{-U -L} lädt den Stand des lokalen Repositorys hoch (z.\,B.\
noch nicht gepushte Commits).

\begin{dshinweis}[title={Zurück heißt nicht: Datenbank zurück}]
\cmd{-{}-zurueck} und das automatische Zurückschalten nehmen Änderungen an der
Datenbank (Migrationen) \emph{nicht} zurück. Bei Updates mit Datenbankänderungen
ist im Notfall die Sicherung vor dem Update der Weg zurück
(Abschnitt~\ref{sec:sicherung}).
\end{dshinweis}

\subsection{Python-Pakete und Django}

Die Versionen aller Pakete stehen fest in \pfad{requirements.txt}. Updates
(z.\,B.\ Sicherheitsupdates von Django) laufen wie jede Änderung: Versionsnummer im
Pull Request erhöhen, Tests abwarten, übernehmen, \cmd{fv-web-update}. Django~6.1
erhält Sicherheitsupdates bis Dezember 2027; der Umstieg auf die Langzeitversion
6.2 (April 2027, Unterstützung bis April 2030) ist einzuplanen.

% ==================================================================
\section{Einstellungen}\label{sec:einstellungen}

Die Einstellungen stehen in \pfad{/etc/fv-web/env} (eine Zeile \cmd{NAME=wert}).
Nach einer Änderung:

\begin{befehl}
sudo systemctl restart fv-web      # startet den Newsletter-Dienst mit neu
\end{befehl}

\begin{center}
\small
\begin{tabularx}{\textwidth}{@{}>{\ttfamily\raggedright}p{52mm}X@{}}
\toprule
\textbf{\sffamily Variable} & \textbf{Bedeutung} \\
\midrule
FV\_SECRET\_KEY            & Geheimer Schlüssel; nicht weitergeben. Ändern meldet alle ab
                             und macht offene Einladungslinks ungültig. \\
FV\_MAIL                   & \cmd{datei} (Mails nur als Dateien in \pfad{/var/lib/fv-web/mails})
                             oder \cmd{smtp} \\
FV\_SMTP\_HOST, \_PORT, \_USER, \_PASSWORD & Zugang zum Mailserver; mit eigenem
                             Mailserver \cmd{127.0.0.1}, Port 25, ohne Benutzer (setzt \cmd{-M}) \\
FV\_SMTP\_TLS              & \cmd{starttls} (Port 587), \cmd{ssl} (Port 465) oder \cmd{aus} \\
FV\_MAIL\_ABSENDER         & Absenderadresse, z.\,B.\ \cmd{webmaster@foerderverein-kf-jf-kreuzau.de};
                             den Namen ergänzt die Webseite \\
FV\_MAIL\_DOMAIN, FV\_MAIL\_SERVER & Domain der Vereinsadressen und Name des Mailservers
                             (setzt \cmd{-M}) \\
FV\_MAIL\_QUOTA\_GB        & Speichergrenze je Postfach in GB (5); nach Änderung Skript mit \cmd{-U} \\
FV\_MAIL\_LOESCHFRIST\_TAGE, \_VORWARNUNG\_TAGE, \_SPERRFRIST\_TAGE & Löschung gesperrter
                             Postfächer (90), Erinnerung vorher (14), Sperre der Adresse (365) \\
FV\_ADMINS                 & Adressen (durch Komma getrennt), die bei Programmfehlern eine
                             Fehlermail mit Details erhalten (nur mit \cmd{smtp}). Die Mails
                             können personenbezogene Daten enthalten: nur die Betreuung eintragen. \\
FV\_ADMINS\_BETREFF        & Betreff-Anfang dieser Fehlermails zum Filtern, z.\,B.\
                             \cmd{"[foerderverein-kf-jf-kreuzau.de]"}. Mit Leerzeichen in
                             Anführungszeichen setzen. \\
FV\_ADMIN\_PFAD            & Pfad der Verwaltung, Standard \cmd{verwaltung/} \\
FV\_NEWSLETTER\_PAUSE      & Sekunden zwischen Probe und Versand an alle (300) \\
FV\_NEWSLETTER\_ABSTAND    & Sekunden zwischen zwei Newsletter-Mails (3) \\
FV\_NEWSLETTER\_MAX\_ANHANG\_MB & Größe, bis zu der Anhänge mitgehen (8) \\
FV\_ZWEI\_FAKTOR\_PFLICHT  & \cmd{1}: Zwei-Faktor für Bearbeitende Pflicht (Standard; nicht abschalten) \\
FV\_LOG\_LEVEL             & Ausführlichkeit des Protokolls, Standard \cmd{INFO} \\
\midrule
FV\_ALLOWED\_HOSTS, FV\_BASIS\_URL, FV\_CSRF\_TRUSTED\_ORIGINS & setzt das Skript aus der
                             Domain -- nicht von Hand ändern \\
FV\_DB\_*, FV\_*\_ROOT, FV\_MAIL\_DATEIEN & setzt das Skript -- nicht ändern \\
\bottomrule
\end{tabularx}
\end{center}

\subsection{Fremden Mailserver eintragen (ohne -M)}

Nur falls die Webseite \emph{ohne} eigenen Mailserver über einen Anbieter verschicken
soll:

\begin{befehl}
sudo nano /etc/fv-web/env
#   FV_MAIL=smtp
#   FV_SMTP_HOST=mail.example.org
#   FV_SMTP_PORT=587
#   FV_SMTP_USER=webmaster@foerderverein-kf-jf-kreuzau.de
#   FV_SMTP_PASSWORD=…
#   FV_SMTP_TLS=starttls
sudo systemctl restart fv-web
\end{befehl}

Danach in der Verwaltung eine Einladung an sich selbst senden und prüfen, ob sie
ankommt. Für die Zustellung müssen SPF, DKIM und DMARC der Domain zum Mailserver
passen.

\subsection{Domain wechseln}

Neue Domain im DNS eintragen und das Skript mit der neuen Domain und \cmd{-U} aufrufen.
Es setzt die domainabhängigen Einstellungen, schreibt Caddy neu und startet die
Webseite neu. Angemeldete Personen müssen sich neu anmelden; bereits verschickte
Einladungslinks zeigen noch auf die alte Adresse. Ein Wechsel der \emph{Mail}-Domain
betrifft alle Vereinsadressen und ist ein eigenes Vorhaben.

% ==================================================================
\section{Datensicherung}\label{sec:sicherung}

\subsection{Was, wann, wohin}

\begin{itemize}
  \item \textbf{Was:} Datenbank (als \cmd{pg\_dump}), alle Anhänge
        (\pfad{/var/lib/fv-web/private}) und die Einstellungen (\pfad{/etc/fv-web}).
        Mit Mailserver zusätzlich alle Postfächer und Archive (\pfad{/var/vmail}),
        DKIM-Schlüssel und Einstellungen (\pfad{/etc/fv-mail}) und das Zertifikat
        (\pfad{/etc/letsencrypt}, wegen des Schlüssels für DANE).
  \item \textbf{Wann:} täglich gegen 03:30 und vor jedem Update; von Hand mit
        \cmd{sudo fv-web-backup}.
  \item \textbf{Wohin:} verschlüsselt mit restic nach \cmd{RESTIC\_REPOSITORY} in
        \pfad{/etc/fv-web/backup.env}. Zusätzlich bleiben die Dumps der letzten 7 Tage
        unverschlüsselt in \pfad{/var/backups/fv-web}.
  \item \textbf{Aufbewahrung:} 30 tägliche und 12 monatliche Stände
        (\cmd{BACKUP\_KEEP\_DAYS}, \cmd{BACKUP\_KEEP\_MONTHS}).
\end{itemize}

\begin{dshinweis}[title={Wichtig: Passwort der Sicherung außerhalb des Servers aufbewahren}]
Ohne \cmd{RESTIC\_PASSWORD} aus \pfad{/etc/fv-web/backup.env} ist die Sicherung
wertlos. Fällt der Server aus, ist auch diese Datei weg. Passwort und Backup-Ziel
deshalb gleich nach der Einrichtung im Passwortsafe des Vereins ablegen:
\begin{befehl}
sudo grep -E '^(RESTIC_PASSWORD|RESTIC_REPOSITORY)=' /etc/fv-web/backup.env
\end{befehl}
\end{dshinweis}

Liegt das Backup-Ziel auf dem Server selbst (Standard \pfad{/var/backups/fv-web-restic}),
schützt es nur vor Bedienfehlern, nicht vor dem Ausfall des Servers. Ein externes Ziel,
z.\,B.\ eine Storage-Box per SFTP, wird mit \cmd{-B} beim Einrichtungsskript gesetzt.
Mit Mailserver wächst die Sicherung mit den Postfächern; Platz beim Ziel einplanen.

\textbf{Überwachung:} Mit \cmd{BACKUP\_PING\_URL} in \pfad{backup.env} meldet die
Sicherung Start, Erfolg und Fehler an einen Überwachungsdienst (z.\,B.\
healthchecks.io), der eine E-Mail schickt, wenn eine Sicherung ausbleibt.

\subsection{Wiederherstellen}

\begin{center}
\small
\begin{tabularx}{\textwidth}{@{}>{\ttfamily\raggedright}p{60mm}X@{}}
\toprule
sudo fv-web-restore -l          & vorhandene Sicherungen anzeigen \\
sudo fv-web-restore -t          & neuesten Stand in eine Test-Datenbank einspielen, Seiten und
                                  Anhänge zählen, wieder löschen -- \textbf{ändert nichts} \\
sudo fv-web-restore -r <stand> -E & Datenbank und Anhänge durch den Stand \textbf{ersetzen}
                                  (Rückfrage; vorher automatische Sicherung) \\
sudo fv-web-restore -r <stand> -d <datei> & einzelne Dateien unter \pfad{/var/lib/fv-web/private}
                                  zurücklegen \\
\bottomrule
\end{tabularx}
\end{center}

\cmd{<stand>} ist die ID aus der Liste oder \cmd{latest}.

\textbf{Einzelnes Postfach zurückholen} (z.\,B.\ versehentlich gelöschte Mails):
Postfach aus der Sicherung in ein Arbeitsverzeichnis holen und mit \cmd{fv-mail
import} einspielen. Vorhandene Mails bleiben dabei erhalten.

\begin{befehl}
sudo bash -c 'set -a; . /etc/fv-web/backup.env; restic -o sftp.args="-i /root/.ssh/fv-web-backup" restore <ID> --host fv-web --include /var/vmail/konten/<nr> --target /root/alt'
sudo fv-mail import /root/alt/var/vmail/konten/<nr>/Maildir <adresse>
\end{befehl}

\subsection{Neuer Server nach Totalausfall}

\begin{enumerate}
  \item Neuen vServer mit Debian 13 bereitstellen, DNS auf die neue Adresse
        (Webseite \emph{und} \cmd{mail.<domain>}, \cmd{mta-sts.<domain>}); PTR beim
        Anbieter neu setzen.
  \item \textbf{Vor} dem Einrichtungsskript das gesicherte Passwort hinterlegen, damit
        das Skript kein neues erzeugt:
\begin{befehl}
ssh root@<neuer-server>
install -d -m 750 /etc/fv-web
( umask 077; echo 'RESTIC_PASSWORD=<aus dem Passwortsafe>' > /etc/fv-web/backup.env )
\end{befehl}
  \item Einrichtungsskript mit der Einstellungsdatei, aber
        vorerst mit \cmd{WITH\_MAIL="0"} und neuer \cmd{SERVER}-Adresse. Den neuen öffentlichen
        Backup-Schlüssel bei der Storage-Box eintragen.
  \item Sicherungen anzeigen und den letzten Stand \textbf{vor dem Ausfall} einspielen:
\begin{befehl}
sudo fv-web-restore -l
sudo fv-web-restore -r <ID vom letzten Tag vor dem Ausfall> -E
\end{befehl}
        \textbf{Nicht} \cmd{latest} verwenden: Das Einrichtungsskript hat auf dem neuen
        Server bereits eine erste Sicherung der noch leeren Datenbank angelegt, und
        \cmd{latest} wäre dieser leere Stand. Die Liste zeigt Datum und Uhrzeit jedes
        Stands.
        Danach \cmd{sudo fv-web-update -f}: installiert die aktuelle Version erneut und
        vergibt dabei die Leserechte für den Mailserver wieder.
  \item Einstellungen aus der Sicherung holen (\pfad{/etc/fv-web} und
        \pfad{/etc/fv-mail} sind darin enthalten) und von Hand in die neue
        \pfad{/etc/fv-web/env} übernehmen; danach \cmd{sudo systemctl restart fv-web}.
        Bei einem SFTP-Ziel:
\begin{befehl}
sudo bash -c 'set -a; . /etc/fv-web/backup.env; restic -o sftp.args="-i /root/.ssh/fv-web-backup" restore <ID> --host fv-web --include /etc/fv-web --include /etc/fv-mail --include /etc/letsencrypt --include /var/vmail --target /root/alt'
sudo cat /root/alt/etc/fv-web/env
\end{befehl}
  \item \textbf{Mailserver:} \pfad{/etc/fv-mail} (DKIM-Schlüssel, Geheimnisse) und
        \pfad{/etc/letsencrypt} aus \pfad{/root/alt} an ihren Platz kopieren, dann
        in der Einstellungsdatei wieder \cmd{WITH\_MAIL="1"} setzen und das Skript mit
        \cmd{-f kreuzau.env -U} aufrufen. So bleiben DKIM- und TLSA-Einträge im DNS
        gültig. Danach die Postfächer zurücklegen:
\begin{befehl}
sudo rsync -a /root/alt/var/vmail/ /var/vmail/
sudo chown -R vmail:vmail /var/vmail
sudo doveadm force-resync -A '*'
sudo fv-mail sync
\end{befehl}
  \item Anmelden und Stichproben prüfen: Seiten, Anhänge, Benutzer, ein Postfach im
        Mailprogramm, \cmd{sudo fv-mail dns}.
\end{enumerate}

% ==================================================================
\section{Überwachung und Störungen}\label{sec:stoerungen}

\subsection{Überwachung von außen: /health und /health/live}\label{sec:health}

Zwei Adressen für einen Überwachungsdienst (z.\,B.\ Uptime Kuma, UptimeRobot,
healthchecks.io), ohne Anmeldung, nie zwischengespeichert -- wie beim Wehrkalender:

\begin{center}
\small
\begin{tabularx}{\textwidth}{@{}>{\ttfamily\raggedright}p{42mm}X@{}}
\toprule
/health/live & \textbf{Lebenszeichen}: Die Anwendung antwortet. Immer
               \cmd{200 \{"status":"ok"\}}, solange die Webseite läuft; keine
               weiteren Prüfungen. \\
/health      & \textbf{Gesamtzustand}: \cmd{200 \{"status":"ok"\}} oder
               \cmd{503 \{"status":"warning", "probleme":[\dots]\}} mit der Liste der
               Probleme. \cmd{"hinweise"} (z.\,B.\ \emph{Neustart nötig}) ändern den
               Status nicht. \\
\bottomrule
\end{tabularx}
\end{center}

\cmd{/health} meldet ein Problem, wenn
\begin{itemize}
  \item die Datenbank nicht antwortet oder der Newsletter-Dienst sich nicht meldet,
  \item ein Dienst nicht läuft (PostgreSQL, Caddy, Webseite, Newsletter; mit
        Mailserver auch Exim, Dovecot, rspamd, Redis, unbound),
  \item die letzte erfolgreiche Datensicherung älter als 26 Stunden ist oder die
        letzte fehlgeschlagen ist,
  \item ein Dateisystem zu 90\,\% oder mehr belegt ist,
  \item ein Zertifikat (Webseite, Mailserver) in weniger als 14 Tagen abläuft oder
        der Mailserver nur ein vorläufiges Zertifikat hat,
  \item mehr als 20 Mails länger als eine Stunde in der Warteschlange warten oder
        der Abgleich bzw.\ die täglichen Aufgaben des Mailservers fehlschlagen,
  \item die Serverprüfung selbst seit mehr als 15 Minuten keine Werte geliefert hat.
\end{itemize}

Die Dienste, Sicherung, Platz und Mailserver prüft \cmd{fv-web-status} als root alle
5~Minuten (\cmd{fv-web-status.timer}) und legt das Ergebnis in
\pfad{/var/lib/fv-web-status/status.json} ab; die Webseite liest es nur. Von Hand:
\cmd{sudo fv-web-status} zeigt das Ergebnis sofort an.

\textbf{Empfehlung für den Überwachungsdienst:} \cmd{/health/live} jede Minute
(Webseite erreichbar?), \cmd{/health} alle 5--10 Minuten mit Benachrichtigung erst
nach zwei aufeinanderfolgenden Fehlern. Die Probleme sind öffentlich lesbar; sie
nennen nur allgemeine Zustände, keine Namen oder Adressen.

\subsection{Zustand prüfen}

\begin{befehl}
sudo fv-web-status                             # Gesamtzustand wie /health
systemctl status fv-web fv-web-newsletter caddy postgresql
systemctl list-timers fv-web-backup.timer      # letzte und nächste Sicherung
curl -s https://foerderverein-kf-jf-kreuzau.de/health   # {"status":"ok"} oder Probleme
df -h /                                        # freier Plattenplatz
sudo fv-web-update --status
sudo fv-mail status                            # Mailserver
sudo exim -bp | head                           # Warteschlange (wartende Mails)
\end{befehl}

\subsection{Protokolle}

\begin{center}
\small
\begin{tabularx}{\textwidth}{@{}>{\ttfamily\raggedright}p{70mm}X@{}}
\toprule
journalctl -u fv-web -e              & Webseite: Fehler, Anfragen \\
journalctl -u fv-web-newsletter -e   & Newsletter-Versand \\
journalctl -u fv-web-backup -e       & Datensicherung \\
journalctl -u caddy -e               & Zertifikat, Webserver \\
sudo less /var/log/fv-web-update.log & alle Updates \\
sudo less /var/log/caddy/fv-web.log  & Zugriffe (JSON) \\
\midrule
sudo exigrep <adresse> /var/log/exim4/mainlog & Weg aller Mails von/an eine Adresse \\
sudo less /var/log/exim4/rejectlog   & abgewiesene Mails mit Grund \\
journalctl -u dovecot -e             & Anmeldungen, Postfächer, Sieve \\
journalctl -u rspamd -e              & Spamprüfung \\
journalctl -u fv-mail-sync -e        & Abgleich mit der Webseite \\
\bottomrule
\end{tabularx}
\end{center}

Die Weboberfläche von rspamd (Statistik, Verlauf) ist nur über einen SSH-Tunnel
erreichbar: \cmd{ssh -L 11334:127.0.0.1:11334 fkjs@<server>}, dann
\url{http://localhost:11334}; Passwort: \cmd{sudo grep RSPAMD /etc/fv-mail/mail.env}.

\subsection{Häufige Störungen}

\begin{description}[style=nextline,leftmargin=0pt,itemsep=0.5\baselineskip]
  \item[Webseite zeigt dauerhaft die Wartungsseite oder lädt nicht]
    \cmd{systemctl status fv-web}, \cmd{journalctl -u fv-web -n 50}. Nach einem
    fehlgeschlagenen Update: \cmd{sudo fv-web-update -{}-status}, ggf.\
    \cmd{-{}-zurueck}. Läuft PostgreSQL?
  \item[Browser meldet Zertifikatsfehler]
    \cmd{journalctl -u caddy -n 50}. Meist zeigt der DNS-Eintrag nicht auf den
    Server oder Port 80 ist gesperrt.
  \item[Newsletter-Statusseite: „Der Versanddienst meldet sich nicht“]
    \cmd{systemctl status fv-web-newsletter}; \cmd{sudo systemctl restart fv-web}.
    Der Versand setzt danach an derselben Stelle fort.
  \item[E-Mails der Webseite kommen nicht an]
    Steht \cmd{FV\_MAIL=smtp} in \pfad{/etc/fv-web/env}? Fehlermeldungen im Protokoll
    von fv-web bzw.\ im Newsletter-Protokoll der Webseite. Mit \cmd{FV\_MAIL=datei}
    liegen die Mails in \pfad{/var/lib/fv-web/mails}. Mit eigenem Mailserver:
    \cmd{sudo exigrep <empfänger> /var/log/exim4/mainlog}.
  \item[Mails nach außen bleiben in der Warteschlange]
    \cmd{sudo exim -bp}; Grund: \cmd{sudo exim -Mvl <id>}. Häufig: Port 25 beim Anbieter
    gesperrt, oder der Empfänger lehnt wegen fehlendem PTR/SPF ab. Erneut versuchen:
    \cmd{sudo exim -qff}. Exim versucht es bis zu 4 Tage lang selbst.
  \item[Mails landen beim Empfänger im Spam]
    \cmd{sudo fv-mail dns} (alles OK?), PTR prüfen, IP-Adresse bei Spamhaus prüfen,
    Testmail an mail-tester.com.
  \item[Eine Mail wurde abgewiesen]
    \cmd{sudo grep <absender> /var/log/exim4/rejectlog}. Gründe: Spam-Punkte zu hoch,
    unzulässiger Anhang (z.\,B.\ \cmd{.exe}, Office mit Makros), Absender auf einer
    Sperrliste, Verteiler nur für Mitglieder, Postfach voll. Greylisting verzögert nur
    („bitte später erneut versuchen“).
  \item[Mailprogramm: Anmeldung schlägt fehl]
    Benutzername ist die vollständige Vereinsadresse, Passwort das der Webseite.
    Passwort einmal auf der Webseite anmelden bzw.\ neu setzen (altes Speicherformat).
    Konto aktiv und Postfach angehakt? Nach 10 Fehlversuchen sperrt fail2ban die
    Adresse für 10 Minuten (bei Wiederholung länger):
    \cmd{sudo fail2ban-client status dovecot}, entsperren mit
    \cmd{sudo fail2ban-client set dovecot unbanip <adresse>} (bzw.\ \cmd{exim}).
  \item[Abwesenheitsnotiz, Archiv oder Belegung aktualisieren sich nicht]
    \cmd{journalctl -u fv-mail-sync -n 30}; von Hand: \cmd{sudo fv-mail sync}.
  \item[Sicherung fehlgeschlagen]
    \cmd{journalctl -u fv-web-backup -n 30}. Häufig: Backup-Schlüssel nicht bei der
    Storage-Box eingetragen oder Ziel voll. Danach \cmd{sudo fv-web-backup} von Hand.
  \item[Platte voll]
    \cmd{df -h}, \cmd{sudo du -sh /var/lib/fv-web/* /var/backups/* /var/log/* /var/vmail/*}.
    Lokale Dumps nach 7 Tagen werden automatisch gelöscht; ein lokales restic-Ziel
    wächst mit. Die Belegung der Postfächer zeigt die Übersicht \emph{Mail} der Webseite.
  \item[Anmeldung per SSH gesperrt]
    fail2ban sperrt eine Adresse nach 5 Fehlversuchen für eine Stunde. Über die
    Konsole des Anbieters: \cmd{sudo fail2ban-client set sshd unbanip <adresse>}.
\end{description}

% ==================================================================
\section{Verwaltungsaufgaben}\label{sec:verwaltung}

Befehle der Anwendung laufen auf dem Server über \cmd{sudo fv-manage <befehl>}.

\subsection{Administrator-Konto der Webseite}

\begin{befehl}
sudo fv-manage vorstand_anlegen <benutzername> <e-mail> --admin
\end{befehl}

Legt das Konto an (oder ergänzt ein vorhandenes), ordnet es der höchsten Stufe mit
Bearbeitungsrecht zu und gibt einen Link zum Setzen des Passworts aus (7 Tage
gültig). Ohne \cmd{-{}-admin} entsteht ein normales Vorstandskonto; mit
\cmd{-{}-stufe <name>} eine andere Stufe.

\subsection{Stufen und Einstellungen der Webseite}

Nur für Administratoren, in der Verwaltung:
\begin{itemize}
  \item \textbf{Stufen}: Name, Rang (höher = mehr Rechte), \emph{darf bearbeiten}.
        Eine Stufe mit zugeordneten Seiten lässt sich nicht löschen.
  \item \textbf{Einstellungen der Webseite}: Name, Slogan, Inhaber (Fußzeile),
        Ansprechpartner, Unterschrift unter E-Mails, weitere Antwortadresse für
        Newsletter.
\end{itemize}

\subsection{Zwei-Faktor-Anmeldung zurücksetzen}

Hat jemand das Smartphone verloren und keine Wiederherstellungscodes:

\begin{befehl}
sudo fv-manage zwei_faktor_zuruecksetzen <benutzername oder e-mail>
\end{befehl}

Vorher sicherstellen, dass die Anfrage wirklich von der Person kommt (z.\,B.\
Rückruf). Personen mit Bearbeitungsrecht müssen die Zwei-Faktor-Anmeldung bei der
nächsten Anmeldung neu einrichten.

\subsection{Gelöschte Seite zurückholen}

\begin{befehl}
sudo fv-manage seite_wiederherstellen          # gelöschte Seiten mit Nummer
sudo fv-manage seite_wiederherstellen 42       # Seite 42 zurückholen
\end{befehl}

Titel, Datum, Sichtbarkeit und Text kommen aus dem Versionsverlauf. Die Dateien der
Anhänge wurden beim Löschen entfernt; der Befehl nennt die fehlenden Dateien und
den passenden \cmd{fv-web-restore}-Befehl. Dafür einen Stand \emph{vor} dem Löschen
wählen (\cmd{sudo fv-web-restore -l}). Danach denselben Befehl noch einmal aufrufen,
er ergänzt die Anhänge.

\subsection{Person ohne E-Mail-Versand einladen}

Solange \cmd{FV\_MAIL=datei} gilt, zeigt die Verwaltung beim Einladen den Link an.
Er muss dann auf anderem Weg an die Person gehen.

\subsection{Mailserver}

Adressen, Postfächer, Aliase, Verteiler und Archive pflegt der Vorstand in der
Webseite (Vorstandsanleitung). Für die Betreuung bleiben:

\begin{center}
\small
\begin{tabularx}{\textwidth}{@{}>{\ttfamily\raggedright}p{70mm}X@{}}
\toprule
sudo fv-mail import <maildir> <adresse> & vorhandene Mails in ein Postfach übernehmen \\
sudo fv-mail sync                       & Abgleich sofort statt nach 5 Minuten \\
sudo exim -Mrm <id>                     & eine Mail aus der Warteschlange löschen \\
sudo doveadm quota get -u <adresse>     & Belegung eines Postfachs \\
sudo doveadm kick <adresse>             & alle Verbindungen einer Person trennen \\
\bottomrule
\end{tabularx}
\end{center}

Ein Postfach wird \emph{nicht} von Hand gelöscht: Konto in der Webseite deaktivieren
bzw.\ Postfach abhaken; Erinnerung und Löschung laufen dann nach den Fristen.

% ==================================================================
\section{Wartungsplan}\label{sec:wartung}

\begin{center}
\small
\begin{tabularx}{\textwidth}{@{}lX@{}}
\toprule
\textbf{Wann} & \textbf{Was} \\
\midrule
monatlich     & Neustart, damit Kernel- und Bibliotheksupdates wirksam werden:
                \cmd{sudo reboot} (Webseite und Mail etwa eine Minute weg; Mails von
                außen werden vom Absender später erneut zugestellt), danach
                \cmd{systemctl status fv-web} und \cmd{sudo fv-mail status}. \\
              & Sicherung prüfen: \cmd{systemctl list-timers fv-web-backup.timer},
                \cmd{journalctl -u fv-web-backup -n 5}. \\
              & Sicherheitsmeldungen zu Django und den Paketen prüfen
                (\url{https://www.djangoproject.com/weblog/}); ggf.\ Update
                (Abschnitt~\ref{sec:update}). \\
              & Mail: Warteschlange leer (\cmd{sudo exim -bpc})? DMARC-Berichte an
                postmaster@ überfliegen. \\
vierteljährlich & Wiederherstellungstest: \cmd{sudo fv-web-restore -t}. \\
              & Benutzerkonten mit dem Vorstand durchsehen: Ausgeschiedene deaktiviert? \\
              & \cmd{sudo fv-mail dns}: alle Einträge OK? IP-Adresse auf Sperrlisten? \\
jährlich      & Zugänge im Passwortsafe prüfen (Backup-Passwort, Storage-Box, GitHub,
                SSH-Schlüssel der Betreuung, SchlundTech, Kundenmenü des Servers). \\
              & Plattenplatz und Größe der Sicherung prüfen (Postfächer wachsen). \\
Frühjahr 2027 & Umstieg auf Django 6.2 LTS. \\
\bottomrule
\end{tabularx}
\end{center}

\end{document}
